\section{Universal Enveloping Algebras}

Write $T(\mathfrak{g} )$ for the tensor algebra on $\mathfrak{g} $. If $\mathfrak{g} $ is a Lie algebra, the universal enveloping algebra can be given as $T(\mathfrak{g})$ modulo the relation that the inclusion $\mathfrak{g} \hookrightarrow T(\mathfrak{g} )$ be a Lie algebra homomorphism and $T(\mathfrak{g} )$ have bracket given by the commutator. It is universal in the sense that precomposition by the inclusion yields
\[
	\Lie_R(-,\Lie )\cong \Alg_R(U,-)
.\]

\begin{proposition}\label{7.3.1}
	If $\mathfrak{g} $ is a Lie algbera, then every $\mathfrak{g} $-module is naturally a $U\mathfrak{g} $-module, and vice versa. This yields natural isomorphisms $\Mod _\mathfrak{g} \cong \Mod _{U\mathfrak{g} } $.
\end{proposition}
\begin{proof}
	Transposing under the tensor-hom adjunction shows that a $\mathfrak{g} $-module is a Lie algebra map $\mathfrak{g} \to \Lie(\End_R(M) )$ for some $R$-module $M$. A $U\mathfrak{g} $-module is an algebra map $U\mathfrak{g} \to \End_R(M) $. The statement now follows by adjointness
	\[
		\Lie_R(\mathfrak{g} ,\Lie (\End_{R} (M) ))\cong \Alg_R(U\mathfrak{g} ,\End_R(M) )
	.\]
\end{proof}
Since the categories of associative algebras have enough projectives and injectives, \cref{7.3.1} shows that so do the categories of $\mathfrak{g} $-modules. The action is simply given by, for $x_1, \ldots ,x_n\in \mathfrak{g} $,
\[
	(x_1 \cdots x_n)m = x_1(\cdots (x_nm)\cdots )
.\]
We then impose bilinearity to cover all of $U\mathfrak{g} $.

\begin{corollary}\label{khsdjfkcnsdjfndks}
	Let $M$ be a $\mathfrak{g} $-module. Then
	\[
		H_{\bullet} (\mathfrak{g} ,M) \cong \operatorname{Tor} _{\bullet} ^{U\mathfrak{g} } (R,M)
	,\]
	\[
		H^{\bullet} (\mathfrak{g} ,M)\cong \operatorname{Ext}_{U\mathfrak{g} } ^{\bullet} (R,M)
	.\]
\end{corollary}
\begin{proof}
	Since we have already observed that $M^\mathfrak{g} \cong \Hom_\mathfrak{g} (R,-) \cong \Hom_{U\mathfrak{g} } (R,-)$, the second statement is immediate by functoriality of derived functors (in the homotopy sense) and homology. For the first, note that the transpose of $0: \mathfrak{g}  \to R$ under $U\dashv \mathrm{Lie} $ is the unique map $U\mathfrak{g} \to R$ mapping $\mathfrak{g}\mapsto 0$. Define the \emph{augmentation ideal} $\mathfrak{J} $ to be the kernel of this map; evidently $\mathfrak{J} $ is the two-sided ideal generated by the image of $\mathfrak{g} $ in $U\mathfrak{g} $. This map is epi, hence the first isomorphism theorem yields $R \cong U\mathfrak{g} /\mathfrak{J} \cong (U\mathfrak{g} )_\mathfrak{g} $. Then
	\[
		R\otimes_{U\mathfrak{g} } M \cong (U\mathfrak{g} /\mathfrak{J} )\otimes _{U\mathfrak{g} } M \cong M / \mathfrak{J} M \cong M / \mathfrak{g} M \cong M_\mathfrak{g} 
	.\]
\end{proof}

\begin{theorem}[Poincar\'e-Birkhoff-Witt]\label{PBW}
	Let $\mathfrak{g} $ be a free $R$-module. Then $U\mathfrak{g} $ is also a free $R$-module. If $\left\{ e_\alpha  \right\} _{\alpha \in A} $ is an ordered basis of $\mathfrak{g} $, then the collection of all elements of the form $e_I$, where $I$ is an increasing sequence, form a basis for $U\mathfrak{g} $.
\end{theorem}
We decline to prove the PBW.
\begin{corollary}\label{lkasj}
	The inclusion $\mathfrak{g} \hookrightarrow U\mathfrak{g} $ is injective, so we may identify $\mathfrak{g} $ with its image.
\end{corollary}
\begin{corollary}\label{askj}
	If $\mathfrak{h} \subseteq \mathfrak{g} $ is a Lie subalgebra, and $R=k$ is a field, then $U\mathfrak{g} $ is a free $U\mathfrak{h} $-module.
\end{corollary}

\begin{exercise}
	Let $M,N$ be left $\mathfrak{g} $-modules. Show that $\Mod _R(M,N)$ is a $\mathfrak{g} $-module by $(xf)(m)=xf(m)-f(xm)$. Then show that $\Mod _\mathfrak{g}(M,N)\cong \Mod _R (M,N)^\mathfrak{g} $.  Show that this induces isomorphisms $\operatorname{Ext} ^{\bullet} _{U\mathfrak{g} } (M,N)\cong H^{\bullet} _{\mathrm{Lie} } (\mathfrak{g} ,\Mod _R(M,N))$. In particular, the global dimension of $U\mathfrak{g} $ equals the Lie algebra cohomological dimension of $\mathfrak{g} $.
\end{exercise}
\begin{proof}
	For the first, we need only show $[x,y]f=x(yf)-y(xf)$. Let $x,y\in \mathfrak{g} $ and $m\in M$. Then since $M,N$ are $\mathfrak{g} $-modules, 
	\begin{align*}
		([x,y]f)(m) &= [x,y]f(m) - f([x,y]m)\\
			    &=xyf(m) - yxf(m) - f(xym - yxm)\\
			    &=xyf(m) - yxf(m) - f(xym) + f(yxm)
	.\end{align*}
	Note that $xf = x\circ f-f\circ x$. Thus we also have
	\[
		x(yf) = x(y\circ f - f\circ y) = x\circ y\circ f - y\circ f\circ x - x \circ f\circ y + f\circ y\circ x
	.\]
	Hence evaluating at $m$ yields
	\[
		(x(yf))(m) = xyf(m) - yf(xm) - xf(ym) + f(yxm)
	.\]
	Similarly,
	\[
		(y(xf))(m) = yxf(m) - xf(ym) - yf(xm) + f(xym)
	.\]
	Thus
	\[
	(x(yf))(m) - (y(xf))(m) \]
	\[
		= xyf(m) - yf(xm) - xf(ym) + f(yxm) - yxf(m) + xf(ym) + yf(xm) - f(xym)
	\]
	\[
		=xyf(m)+f(yxm) - yxf(m) - f(xym)
	.\]
	They agree, hence $\Mod _R(M,N)$ is a $\mathfrak{g} $-module in the way claimed. 

	Now let $f$ be a $\mathfrak{g} $-invariant of $\Mod _R(M,N)$. Then $0=xf = x\circ f - f\circ x$ for all $x\in \mathfrak{g} $, and thus $x\circ f = f\circ x$. Thus $f$ is $\mathfrak{g} $-linear. Conversely if $f$ is $\mathfrak{g} $-linear, then $f\circ x = x\circ f$, and thus $xf = x\circ f-f\circ x =0$, so $f$ is a $\mathfrak{g} $-invariant of $\Mod_R(M,N)$.
\end{proof}

